% Independent proof audit. Fragment for the main manuscript.
% Requires amsmath, amssymb, and a theorem environment supplied by the parent.
\section{Two forced-symbol identifications and the 246-state machine}
\label{sec:forced246}

All state names in this section are the original numerical labels, with
254 the separate terminal label.  Let $M_{250}$ be the prefix-carry and
pulse machine in the preceding construction.  Its legal runtime frames
have nine unary primary blocks and a nonempty finite spare suffix.  Define
$M_{246}$ by deleting states $0,1,3,97$, changing just the two surviving rows
\[
 \delta(60,0)=(0,R,20),\qquad
 \delta(255,0)=(1,L,255),
\]
and taking state 95 as the blank-tape initial state.  Every other surviving
row is unchanged.  This is a total two-symbol machine with 246 working
states, 492 transition entries, and one terminal transition.  The terminal
entry is still $\delta(77,1)=(1,R,254)$, including its final write and move.

The reduction uses information already guaranteed by the tape format and
by a preceding write.  It preserves the complete original proof-search
program, including all eight logical rows, the equality machinery, both
freshness tests, and the CP/GIC bundle.  Its correctness does not depend on
the rejected equality-free compiler experiments.

\subsection{The marking-state identification}

Write $B(v)=0\,1^{v+1}$.  The old marking state has rows
\[
 \delta_{250}(60,0)=(0,R,0),\qquad
 \delta_{250}(0,1)=(0,L,2).
\]
The existing state 20 has
\[
 \delta_{250}(20,0)=(1,L,2),\qquad
 \delta_{250}(20,1)=(0,L,2).
\]
Thus its read-one action agrees exactly with state 0's read-one action.
At a legal marking call, state 60 reads the zero separator of the selected
unmarked block $B(v)$.  Its right neighbor is one for every $v\geq0$,
including the zero-valued register.  The two executions are therefore
\[
 \begin{array}{rcl}
 M_{250}:&60\xrightarrow{0}0\xrightarrow{1}2,\\
 M_{246}:&60\xrightarrow{0}20\xrightarrow{1}2.
 \end{array}
\]
They have the same duration, head position, tape writes, and endpoint.
Neither execution uses the read-zero row of the intermediate state.
The remainder of the marking return scan is literally unchanged.

This argument covers every physical register selector.  If the selected
register is preceded by the block list $A$, the selector scan reaches its
separator after $w(A)=\sum_{v\in A}(v+2)$ steps.  The return scan has the
same width.  The old marking theorem therefore retains its exact duration
$2w(A)+3$, its erased first unary one, and its endpoint state 2 for the
first block or state 6 for the other eight blocks.  Both exterior tape
contexts are retained.

The original transition table has only one incoming edge to state 0,
namely the changed $60/0$ edge.  After redirecting it, states 0, 1, and 3
have no incoming edge from the retained control.  The only incoming edges
to state 1 were $0/0$, $1/1$, and $3/1$; the only incoming edge to state 3
was $1/0$.  Hence all three can be deleted together.  This is an actual
three-state deletion.  The known symbol at marking is what licenses the
redirection; ordinary exact state equivalence would not license it.

\subsection{The pulse-state identification}

The relevant prefix250 rows are
\[
\begin{array}{c|cc}
 &0&1\\\hline
75 &(1,R,255)&(1,L,63)\\
255&(1,L,97)&(1,L,131)\\
97 &(1,R,83)&(1,L,131).
\end{array}
\]
The read-one rows of states 97 and 255 coincide.  Every recurrent entry
to state 97 follows $255/0$, which itself follows $75/0$ in a local pulse.
That preceding transition has already written one at the cell to which
$255/0$ returns.  The returning state 97 consequently reads one.

More explicitly, start in state 75 at coordinate $x$ reading zero.
If the right neighbor is one, both machines reach state 131 at coordinate
$x$ after two transitions, with the same whole tape.  If that neighbor is
zero, the three-step state sequences are
\[
 \begin{array}{rcl}
 M_{250}:&75\longrightarrow255\longrightarrow97\longrightarrow131,\\
 M_{246}:&75\longrightarrow255\longrightarrow255\longrightarrow131.
 \end{array}
\]
The first transition writes one at $x$, the second writes one at $x+1$
and moves left, and the third reads the deposited one at $x$ and moves
left again.  Both executions reach state 131 at $x-1$ after exactly three
steps and agree at every tape coordinate.  This proves the pulse identity
for either possible neighbor bit, with arbitrary exterior contexts.

The only old incoming edge to 97 is $255/0$.  Redirecting it to 255 removes
97 completely.  The old blank initialization at $97/0$ is consequently
discarded; a separately checked blank run from state 95 provides the new
initializer.  No assertion is made that the old and new blank prefixes
agree.

\subsection{Direct checks of the universal local identities}

An independent constructor in \texttt{audit246.py} applies precisely these
edits to the separately reconstructed prefix250 table.  The check is not
limited to the two informal state diagrams.  The inherited mathematical
induction checker is rerun against the new table for:
\begin{itemize}
\item the four arbitrary-length block identities for the seeker, carry,
return, and rollback scans;
\item the positive-decrement, zero-decrement, and increment macros, with
symbolic prefix and suffix lists and arbitrary natural register values;
\item all nine marking entries, with a separate arbitrary natural length
for every primary register.
\end{itemize}
All sixteen checks pass.  The proof rules are literal transition steps,
induction on unary self-loop lengths, structural induction on finite lists
of register blocks, and the finite-word identity $C(A)1=1B(A)$ for
$C(v)=10\,1^v$.  Opaque exterior contexts are never inspected.  The
zero-decrement theorem retains its nonempty following-block premise;
no zero-spare or bounded-register assumption is introduced.

The independent run then reparses all 6,459 nodes of the literal
\texttt{Prog254.v} decoder and checks every one of its 6,460 leaves,
including both guards at each decrement.  All 8,391 instruction cases
pass against the changed table.  They are parameterized by every extra
gap $n\geq0$, every primary register vector satisfying the named guard,
and every nonempty finite spare list.  The guard names the same physical
register as the selected block.  Each accepted case contains at least
13 literal steps, and no case requires more than 497 symbolic rules.

These are explicit finite symbolic-induction certificates and ordinary
mathematical proofs of their rule meanings.  They are not a new Lean or
Rocq build of the 246-state endpoint.  A proof-assistant port can retain
the original logical and model-theoretic modules and specialize the tape,
bootstrap, and runtime modules to this new table.

\section{Blank initialization and the inherited consistency calibration}
\label{sec:calibration246}

Two separately written literal interpreters, one using a set of one-cell
coordinates and one using a bounded dense array with checked bounds,
execute the complete 410,735-transition blank prefix.  They agree on the
entire endpoint:
\[
 q=96,\qquad h=-5,\qquad \text{header}=176,\qquad d=181,\qquad P=[].
\]
The unique left marker is at $-6$.  The physical primary values are
\[
 (0,0,0,3,0,0,3,0,0),
\]
followed by 167 spare blocks of value one.  There are exactly 351 one-cells.
Every remaining tape cell is zero.  The recorded finite support is
reconstructed from this block description and compared against both
interpreters; a finite window is not substituted for a whole-tape claim.

The physical positions of logical registers $(N,T,W,A,D,H,X,Y,L)$ remain
\[
 (6,2,0,4,3,7,8,5,1).
\]
Thus the actual logical root vector is
\[
 R_*=(3,0,0,0,3,0,0,0,0).
\]
The nonzero fields are the next-candidate counter $N=3$ and numerical
parameter $D=3$.  Current formula, cached formula, and formula stack are
all zero.  The inherited theorem
\texttt{DirtyProgram254.six\_nonformula\_path\_iff}
applies with arguments $(N,A,D,X,Y,L)=(3,0,3,0,0,0)$.
Its safety premises are the two instances of
\texttt{Code254.C254.truth} and \texttt{Enum254.good\_zero}.
Starting the candidate counter at three loses no possible refutation:
\texttt{EnumSafeStart254.complete\_from} supplies arbitrarily large padded
encodings of any refutation and proves the required scheduler coverage.

\subsection{Unbounded runtime composition}

The complete recurrent frame has state 96 and tape
\[
 0^{\mathbb Z_{<h-1}}\;1\;P\;0^{d-|P|}\;1\;
 B(V(R)S)\;0^\omega,
\]
where $d\geq40$, $|P|\leq d$, $S$ is a nonempty finite spare list, and
$P$ satisfies the inherited clean-path predicate.  The lower bound concerns
the total head-to-header distance, not just the zero gap after the current
path.  The actual bootstrap has $181\geq40$ and an empty, hence clean,
path.

Each nonterminal logical instruction has a positive finite binary segment
to the next complete recurrent frame.  The clean-path condition rebases
the actual zero-padded PC to its consumed leaf path.  Positive decrement
increases the header distance by one; a carry beyond the left marker also
increases that distance by one.  Other instructions preserve the distance.
Increment appends a value-one spare, and every other instruction preserves
the spare suffix.  Consequently neither the margin nor the nonempty spare
resource can be exhausted during an arbitrarily long computation.

The two forced-symbol identifications preserve every local contract,
including its whole-tape endpoint and exact duration.  The full instruction
certificate therefore establishes the same positive-segment simulation
for $M_{246}$ as for the prefix250 runtime.  Finite induction composes these
segments along every finite logical computation.

Halting reflection requires no time bound and no infinite choice of
segments.  Suppose the binary machine halts in $t$ steps from a related
frame.  If the logical instruction is terminal, decoder coverage provides
the logical halt.  Otherwise its positive segment of length $k$ ends in an
active related frame.  Absorption of terminal configurations forces $k<t$;
the remaining binary halt has time $t-k<t$.  Strong induction on $t$ then
produces a finite logical halt.  Thus the relation gives eventual halting
equivalence in both directions.

Combining the actual blank prefix with this equivalence and the dirty-root
theorem yields
\[
 \operatorname{Halt}(M_{246},\mathrm{blank})
 \quad\longleftrightarrow\quad C_{254}(1).
\]
The literal decoder-to-structured-program adapter is
\texttt{Bisim254.prog254\_halts\_iff}; the numerical proof-system
statement is exactly the original one, not a replacement calculus.

\subsection{What is proved in ZFC}

Fix the ordinary syntactic consistency statement
$\operatorname{Con}(\mathrm{ZFC})=\neg\operatorname{Ref}_{\mathrm{ZFC}}$,
where refutations are finite proofs in a fixed recursive first-order
presentation.  The inherited endpoint chain is
\[
 C_{254}(1)
 \ \longleftrightarrow\ 
 \operatorname{Ref}_{\mathrm{NativeAx}}
 \ \longleftrightarrow\ 
 \operatorname{Ref}_{\mathrm{ZF}}
 \ \longleftrightarrow\ 
 \operatorname{Ref}_{\mathrm{ZFC}}.
\]
Here \texttt{NativeAx} means precisely CP, all first-order-formula GIC
instances, and raw irreflexivity.  The relevant inherited declarations are
\texttt{NativeTheory254.c254\_refutation\_iff},
\texttt{WFTheory.native\_refutation\_iff\_zf}, and
\texttt{ChoiceCalibration254.zf\_refutation\_iff\_zfc}.

The model-theoretic proof is legitimate inside ZFC.  Finite syntax, proofs,
and tape configurations are coded by natural numbers or finite sequences.
First-order completeness supplies a set-coded model from consistency; it
does not supply a transitive or externally well-founded model.  The native
CP/GIC model gives a set-sized bisimulation quotient, and the internally
well-founded predicate supplies the Foundation subdomain used by
\texttt{WFModel.quotient\_models\_ZF}.  Conversely, every model of the
specified ZF presentation, which includes Foundation, satisfies the native
axioms by the internal rank and Collection constructions.  The ZFC
comparison forms the internal constructible subdomain of any set model of
ZF, proving Choice there.  Although that model regards its $L$ as a class,
its external domain is a set, so the resulting subdomain is a set in the
metatheory.  Satisfaction is satisfaction for these set structures, not a
truth predicate for the universe.

All machine arguments above are finite combinatorial inductions, and all
model constructions and satisfaction recursions are ordinary ZFC
constructions.  Accordingly their composition gives the conventional
mathematical theorem
\[
 \mathrm{ZFC}\ \vdash\
 \bigl(\neg\operatorname{Halt}(M_{246},\mathrm{blank})
       \longleftrightarrow\operatorname{Con}(\mathrm{ZFC})\bigr).
\]
This assertion does not assume that ZFC is consistent.  A consistency
assumption is used only when drawing the separate consequence that the
machine really does not halt and that ZFC cannot prove its nonhalting.
The existing 252-state formal endpoint is an external Lean metatheorem;
the present proof supplies an ordinary ZFC derivation for the new table,
not an asserted PA-internal proof object or a freshly compiled 246-state
Lean/Rocq endpoint.
